\section{Axiomatic Framework}
    \label{sec:axioms}
    
    \subsection{Preliminaries}
        
        We formalize Swirl-String Theory (SST) as an axiomatic continuum theory defined over a three-dimensional spatial domain $\mathcal{M} \subset \mathbb{R}^3$ equipped with a preferred time parameter $t \in \mathbb{R}$.
        
        The theory is constructed from a minimal primitive set $\mathcal{P}$ (Section~\ref{sec:foundations}) and a set of axioms $\mathcal{A}$ governing admissible configurations and dynamics.
    
    \subsection{Axiom 1: Continuum Substrate}
        
        \textbf{Statement.}
        There exists a continuous, incompressible, inviscid medium filling $\mathcal{M}$, characterized by a constant background density $\rho_f$.
        
        \begin{align}
            \nabla \cdot \mathbf{v} = 0
        \end{align}
        
        \textbf{[ORTHODOX]}: This corresponds to the standard incompressible Euler framework.
        
        \textbf{Interpretation.}
        All physical structures arise as configurations within this medium; no discrete particles are assumed at the fundamental level.
    
    \subsection{Axiom 2: Existence of Stable Vortex Filaments}
        
        \textbf{Statement.}
        The medium admits stable, localized, topologically nontrivial vortex filament solutions.
        
        \textbf{[ORTHODOX PART]}: Vortex filaments are known solutions in fluid dynamics.
        
        \textbf{[SPECULATIVE EXTENSION]}: Such filaments are assumed to be dynamically stable and persist as fundamental structures.
        
        \textbf{Mathematical form.}
        A filament is represented as a closed embedding:
        
        \begin{align}
            X: S^1 \to \mathcal{M}, \quad \sigma \mapsto X(t,\sigma)
        \end{align}
    
    \subsection{Axiom 3: Circulation Quantization}
        
        \textbf{Statement.}
        Circulation around any stable filament is quantized:
        
        \begin{align}
            \Gamma = \oint \mathbf{v} \cdot d\ell = n \Gamma_0, \quad n \in \mathbb{Z}
        \end{align}
        
        \textbf{[ORTHODOX]}: Quantized circulation appears in superfluid systems.
        
        \textbf{[SPECULATIVE]}: Elevated here to a universal invariant.
    
    \subsection{Axiom 4: Compton Anchoring}
        
        \textbf{Statement.}
        The characteristic angular frequency of the vortex core satisfies:
        
        \begin{align}
            \frac{v_{\swirlarrow}}{r_c} = \omega_c
        \end{align}
        
        \textbf{[ORTHODOX]}: $\omega_c$ is the Compton frequency.
        
        \textbf{[SPECULATIVE]}: Identified as the intrinsic rotational frequency of the vortex core.
        
        \textbf{Consequence.}
        The core radius is not independent but derived (Eq.~\ref{eq:rc_definition}).
    
    \subsection{Axiom 5: Delay as a Constitutive Principle}
        
        \textbf{Statement.}
        Finite propagation time along closed vortex structures is a fundamental dynamical ingredient.
        
        \begin{align}
            \tau = \frac{L}{v_{\swirlarrow}}
        \end{align}
        
        \textbf{[DERIVED]}: Follows from finite signal speed.
        
        \textbf{[AXIOMATIC ROLE]}: The delay is not negligible but structurally relevant.
        
        \textbf{Interpretation.}
        Temporal nonlocality is intrinsic to the dynamics.
    
    \subsection{Axiom 6: Energy Localization}
        
        \textbf{Statement.}
        Energy is localized in vortex structures and determines mass via:
        
        \begin{align}
            E = m c^2
        \end{align}
        
        \textbf{[ORTHODOX]}: Relativistic energy relation.
        
        \textbf{[SPECULATIVE]}: Interpreted as rotational energy of the vortex configuration.
    
    \subsection{Axiom 7: Swirl-Clock Relational Time}
        
        \textbf{Statement.}
        Local time is encoded in a scalar field derived from fluid kinematics:
        
        \begin{align}
            S_{(t)} = \sqrt{1 - \frac{v^2}{c^2}}
        \end{align}
        
        \textbf{[ORTHODOX]}: Matches Lorentz time dilation.
        
        \textbf{[DERIVED]}: Emerges from local velocity field.
        
        \textbf{Interpretation.}
        Time is not fundamental but relational to the state of motion in the medium.
    
    \subsection{Axiom 8: Density Saturation and Core Structure}
        
        \textbf{Statement.}
        The vortex core exhibits a saturated density $\rho_{core}$ determined by mechanical constraints of the medium.
        
        \textbf{[DERIVED]}: From force balance and geometric closure (Eq.~\ref{eq:core_density}).
        
        \textbf{Interpretation.}
        Matter corresponds to regions where energy density reaches a structural limit.
    
    \subsection{Logical Structure of the Axioms}
        
        The axioms form a dependency hierarchy:
        
        \begin{align}
            \text{Continuum}
            \rightarrow \text{Vortex Structures}
            \rightarrow \text{Quantization}
            \rightarrow \text{Delay Dynamics}
            \rightarrow \text{Mass/Energy Emergence}
        \end{align}
    
    \subsection{Minimality Principle}
        
        \textbf{Statement.}
        The axioms are minimal in the sense that removing any one of them destroys at least one of the following:
        
        \begin{itemize}
            \item existence of stable structures
            \item emergence of discrete spectra
            \item compatibility with known physics
        \end{itemize}
    
    \subsection{Philosophical Interpretation}
        
        The axiomatic structure implies:
        
        \begin{itemize}
            \item Particles are not fundamental objects but stable dynamical configurations
            \item Discreteness is emergent, not postulated
            \item Time is relational, not absolute
        \end{itemize}
    
    \subsection{Consistency Condition}
        
        A valid SST realization must satisfy:
        
        \begin{itemize}
            \item Compatibility with orthodox physics in appropriate limits
            \item Internal logical closure
            \item Explicit traceability from axioms to predictions
        \end{itemize}
    
    \subsection{Canonical Statement}
        
        \begin{center}
            \textit{
                Swirl-String Theory is defined as the minimal continuum theory in which stable vortex structures, governed by quantized circulation and delay dynamics, give rise to discrete physical phenomena consistent with known physics.
            }
        \end{center}